\documentclass{beamer}
\usepackage{listings}
\usepackage{amsmath}
\usepackage{caption}
\usepackage{fancyvrb}
\usepackage{xcolor}
\usepackage{eso-pic}

% Reproducible builds
\pdfinfoomitdate=1
\pdftrailerid{}
\pdfsuppressptexinfo=-1

\usepackage{color}
\definecolor{bluekeywords}{rgb}{0.13,0.13,1}
\definecolor{greencomments}{rgb}{0,0.5,0}
\definecolor{turqusnumbers}{rgb}{0.17,0.57,0.69}
\definecolor{redstrings}{rgb}{0.5,0,0}
\definecolor{doge}{RGB}{240,218,149}
\definecolor{lightdoge}{RGB}{245,228,169}
\definecolor{grey}{RGB}{186,171,169}

\setbeamercolor{background canvas}{bg=doge}
\setbeamercolor{frametitle}{bg=lightdoge,fg=black}
\setbeamercolor{section number projected}{bg=doge,fg=black}
\setbeamercolor{section in toc shaded}{bg=doge, fg=grey}
\setbeamercolor{subsection in toc shaded}{bg=doge, fg=grey}
\useinnertheme{rectangles}
\useoutertheme{shadow}

\newcommand\AtPagemyUpperLeft[1]{\AtPageLowerLeft{%
\put(\LenToUnit{0.9\paperwidth},\LenToUnit{0.88\paperheight}){#1}}}
\AddToShipoutPictureFG{
  \AtPagemyUpperLeft{{\includegraphics[width=1cm,keepaspectratio]{BabyDoge.png}}}
}%

\lstdefinelanguage{FSharp}
    {morekeywords={val, let, new, match, with, rec, open, module, namespace, type, of, member, and, for, in, do, begin, end, fun, function, try, mutable, if, then, else},
     keywordstyle=\color{bluekeywords},
     sensitive=false,
     morecomment=[l][\color{greencomments}]{///},
     morecomment=[l][\color{greencomments}]{//},
     morecomment=[s][\color{greencomments}]{{(*}{*)}},
     morestring=[b]",
     stringstyle=\color{redstrings}
     }

\title{Property-Based Testing}
\author{Patrick Stevens}
\institute{G-Research}
\date{Doge Conf 2019}

\lstnewenvironment{fslisting}
  {
    \lstset{
        language=FSharp,
        basicstyle=\small\ttfamily,
        breaklines=true,
        columns=fullflexible}
  }
  {
  }

\begin{document}

\begin{frame}
\titlepage
\end{frame}

\section{What's the problem?}
\subsection{A program to test}

\begin{frame}
\tableofcontents
\end{frame}

\begin{frame}
\tableofcontents[currentsection]
\end{frame}

\begin{frame}
\frametitle{Let's have something to test}
Interval set: a space-efficient set of integers

Defining example:
$$\{{\color{red}1,2,3,4},{\color{cyan}8,9,10}\}$$ $$[{\color{red}(1, 4)}, {\color{cyan}(8, 10)}]$$
\end{frame}

\begin{frame}[fragile]
\frametitle{Public API}
\begin{fslisting}
type IntervalSet

[<RequireQualifiedAccess>]
module IntervalSet =
    val empty : IntervalSet
    val add : int -> IntervalSet -> IntervalSet
    val contains : int -> IntervalSet -> bool
\end{fslisting}

\end{frame}

\begin{frame}
\frametitle{Implementation}
This is a presentation about testing.

Yes, there will be a bug somewhere.
\end{frame}

\begin{frame}[fragile]
\frametitle{Data structure}

\begin{fslisting}
type private Interval =
    {
        Min : int
        Max : int
    }

type IntervalSet = private IntervalSet of Interval list

\end{fslisting}
\end{frame}

\begin{frame}[fragile]
\frametitle{Implementation: empty}

\begin{fslisting}
[<RequireQualifiedAccess>]
module IntervalSet =
    let empty = IntervalSet []

\end{fslisting}
\end{frame}

\begin{frame}[fragile]
\frametitle{Implementation: insertion}
\begin{fslisting}
let private rec add' (a : int) (ls : Interval list) =
    match ls with
    | [] -> [{ Min = a ; Max = a }]
    | _ -> ...
\end{fslisting}
\end{frame}

\begin{frame}[fragile]
\frametitle{Implementation: insertion}
\begin{fslisting}
let private rec add' (a : int) (ls : Interval list) =
    match ls with
    | [] -> ...
    | interval :: is ->
        if interval.Min <= a && a <= interval.Max then
            ls // no need to add, it's already there
        else ...
\end{fslisting}
\end{frame}

\begin{frame}[fragile]
\frametitle{Implementation: insertion}
\begin{fslisting}
let private rec add' (a : int) (ls : Interval list) =
    match ls with
    | [] -> ...
    | interval :: is ->
        if // already there...
        elif interval.Min - 1 = a then
            { interval with Min = interval.Min - 1 }
            :: is // augment this interval to contain a
        elif interval.Max + 1 = a then
            { interval with Max = interval.Max + 1 }
            :: is // augment this interval to contain a
        ...
\end{fslisting}
\end{frame}

\begin{frame}[fragile]
\frametitle{Implementation: insertion}
\begin{fslisting}
let private rec add' (a : int) (ls : Interval list) =
    match ls with
    | [] -> ...
    | interval :: is ->
        if // already there...
        elif // can be added to this interval...
        else // can't add it here; recurse
            add' a is
\end{fslisting}
\end{frame}

\begin{frame}[fragile]
\frametitle{Implementation: insertion}
\begin{fslisting}
[<RequireQualifiedAccess>]
module IntervalSet =
    let add (a : int) (IntervalSet intervals) =
        add' a intervals
        |> IntervalSet
\end{fslisting}
\end{frame}

\begin{frame}[fragile]
\frametitle{Implementation: containment}

\begin{fslisting}
let rec contains (a : int) (IntervalSet ls) =
    ls
    |> List.tryFind (fun interval ->
        interval.Min <= a && a <= interval.Max)
    |> Option.isSome
\end{fslisting}
\end{frame}

\subsection{Testing the program}

\begin{frame}[fragile]
\frametitle{Start testing!}
Helper function for tests:
\begin{fslisting}
let create (is : int list) : IntervalSet =
    is
    |> List.fold
        (fun set i -> IntervalSet.add i set)
        IntervalSet.empty
\end{fslisting}
\end{frame}
\begin{frame}
\frametitle{What can we test?}
We should test some different lists and their resulting IntervalSets.

\begin{itemize}
    \item $[3; 4]$ contains $5$? (No.)
    \item $[3; 5]$ contains $5$? (Yes.)
    \item $[3; 4; 5]$ contains $4$? (Yes.)
\end{itemize}
\end{frame}

\begin{frame}[fragile]
\frametitle{The test cases}
\begin{fslisting}
    create [3; 4]
    |> IntervalSet.contains 5
    |> shouldEqual false

    create [3; 5]
    |> IntervalSet.contains 5
    |> shouldEqual true

    create [3; 4; 5]
    |> IntervalSet.contains 4
    |> shouldEqual true
\end{fslisting}
\pause
Hooray, the tests pass!
\end{frame}

\subsection{But can I really trust myself?}

\begin{frame}
\frametitle{... but is it right?}
\end{frame}

\begin{frame}
\frametitle{... but is it right?}

I don't know!

\begin{itemize}
\item I'm lazy
\item I'm stupid
\item I hate testing
\end{itemize}

Isn't there a better way?
\end{frame}

\begin{frame}[fragile]
\frametitle{FsCheck can help!}

Sneak peek: FsCheck will tell us that this implementation is wrong.

\begin{Verbatim}[commandchars=\\\{\}]
\textcolor{red}{Falsifiable}, after 35 tests (15 shrinks)
  (StGen 9514417537,296661223)):
Original:
[0; 0; 0; 0; 0; 0; 0; 12; 1; -2; 0; 0; 0; 0; 0; 0; 0]
12
Shrunk:
[12; 0]
12
\end{Verbatim}

\end{frame}

\begin{frame}[fragile]
\frametitle{FsCheck's test}
The last two lines FsCheck gave us were:
\begin{verbatim}
[12; 0]
12
\end{verbatim}

\pause

FsCheck found this test:

\begin{fslisting}
create [12; 0]
|> IntervalSet.contains 12
|> shouldEqual true
\end{fslisting}
\end{frame}

\section{Introduction to FsCheck}
\subsection{FsCheck's view of the world}

\begin{frame}
\tableofcontents[currentsection]
\end{frame}

\begin{frame}
\frametitle{Why do you test?}

\begin{itemize}
    \item Your program does what you want it to.
    \item Your program doesn't do what you don't want it to.
\end{itemize}
\end{frame}

\begin{frame}
\frametitle{How do you normally test?}

\begin{enumerate}
    \item Come up with examples.
    \item Work out what your program should do on those examples.
    \item Run the program and check it did what you wanted.
\end{enumerate}
\end{frame}

\begin{frame}
\frametitle{But what are you \emph{really} doing?}

You're testing \emph{properties} through \emph{representative examples}.
\end{frame}

\begin{frame}
    Why not just \emph{test properties}?
\end{frame}

\begin{frame}
\frametitle{FsCheck tests properties automatically}

\begin{itemize}
    \item Find edge cases
    \item Find large, complicated cases
    \item Shrink large cases automatically
    \item Make sure you can repeat any failures (the TDD way!)
\end{itemize}
\end{frame}

\subsection{Back to the example}

\begin{frame}
\frametitle{The failing property for IntervalSet}
\begin{itemize}
\item Create an IntervalSet from a list of integers\dots
\item then check for containment\dots
\item should be the same as checking the original list.
\end{itemize}
\end{frame}

\begin{frame}[fragile]
\frametitle{The failing property for IntervalSet, in code}
\begin{fslisting}
let property (ints : int list) (toCheck : int) : bool =
    create ints
    |> IntervalSet.contains toCheck
    |> (=) (List.contains doesContain ints)
\end{fslisting}

% Remember, FsCheck told us that this property failed.
\end{frame}

\begin{frame}[fragile]
\frametitle{Invoke FsCheck}
\begin{fslisting}
open FsCheck

let property (ints : int list) (toCheck : int) : bool =
    create ints
    |> IntervalSet.contains toCheck
    |> (=) (List.contains doesContain ints)

[<Test>]
let testProperty () =
    Check.QuickThrowOnFailure property
\end{fslisting}

% Remember, FsCheck told us that this failed.
\end{frame}

\begin{frame}[fragile]
\frametitle{FsCheck's output, revisited}
\begin{Verbatim}[commandchars=\\\{\}]
Falsifiable, after \textcolor{red}{35} tests (15 shrinks)
  (StGen 9514417537,296661223)):
Original:
[0; 0; 0; 0; 0; 0; 0; 12; 1; -2; 0; 0; 0; 0; 0; 0; 0]
12
Shrunk:
\textcolor{cyan}{[12; 0]}
\textcolor{cyan}{12}
\end{Verbatim}

\hfill \break

FsCheck constructed \textcolor{red}{35} tests before finding a failure.
It then shrank the test case to \textcolor{cyan}{the smallest failure} it could find.
\end{frame}

\subsection{Advantages}

\begin{frame}
\frametitle{Advantages}

\begin{itemize}
\item No thought required!
\item Perfectly reproducible
\item Edge cases automatically examined closely
\item Randomised testing increases coverage
\end{itemize}

\hfill \break

(Make sure you explicitly test any failures FsCheck finds, so that nothing is lost to the mists of time!)
\end{frame}

\subsection{What was the bug?}
\begin{frame}[fragile]
\frametitle{The bug}
\begin{fslisting}
let private rec add' (a : int) (ls : Interval list) =
    match ls with
    | [] -> [{ Min = a ; Max = a }]
    | interval :: is ->
        if (* contains *) then
            ls
        elif (* is 1 below *) then
        elif (* is 1 above *) then
        else
            add' a is // <-- Oh no!
\end{fslisting}
\end{frame}

\begin{frame}[fragile]
\frametitle{The fix}
\begin{fslisting}
let private rec add' (a : int) (ls : Interval list) =
    match ls with
    | [] -> [{ Min = a ; Max = a }]
    | interval :: is ->
        if (* contains *) then
            ls
        elif (* is 1 below *) then
        elif (* is 1 above *) then
        else
            interval :: add' a is
\end{fslisting}
\end{frame}

\begin{frame}[fragile]
\frametitle{FsCheck is happy}
\begin{verbatim}
Ok, passed 100 tests.
\end{verbatim}
\end{frame}

\subsection{Serialisers}
\begin{frame}
\frametitle{Serialisers: the hello-world of black-box testing}

A serialiser is defined by a property:

\hfill \break

\noindent\fbox{%
Writing an object, then reading it in, gives the original object.
}

\end{frame}

\begin{frame}[fragile]
\frametitle{Signature of a serialiser}

\begin{fslisting}
[<RequireQualifiedAccess>]
module FancyThing =
    val toString : FancyThing -> string
    val parse : string -> FancyThing option
\end{fslisting}
\end{frame}

\begin{frame}[fragile]
\frametitle{Test the serialiser}

\begin{fslisting}
[<Test>]
let roundTripTest () =
    let property (t : FancyThing) : bool =
        t
        |> FancyThing.toString
        |> FancyThing.parse
        |> (=) (Some t)
    Check.QuickThrowOnFailure property
\end{fslisting}
\end{frame}

\begin{frame}
\frametitle{What FsCheck gave us}
Without needing to know anything about the implementation, FsCheck was still able to produce useful tests!
\end{frame}

\section{Metatesting}
\tableofcontents[currentsection]

\begin{frame}
\frametitle{Metatesting}
A technique you should use to make your testing more effective.
\end{frame}

\subsection{Was the testing comprehensive?}

\begin{frame}[fragile]
\frametitle{How can we be sure we tested enough?}

Recall the property:
\begin{fslisting}
let property (ints : int list) (toCheck : int) : bool =
    create ints
    |> IntervalSet.contains toCheck
    |> (=) (List.contains doesContain ints)

\end{fslisting}

\pause
By fluke (or my incompetence), FsCheck might never generate a ``yes, does contain'' case.
\end{frame}

\begin{frame}
\frametitle{Gather metrics}
F\# is impure and side-effectful, so it's extremely easy to gather metrics.

\hfill \break

(Instrumentation is an example where purity really does make things harder!)
\end{frame}

\begin{frame}[fragile]
\frametitle{Gather metrics in the property}
\begin{fslisting}
let property
    (positives : int ref) (negatives : int ref)
    (ints : int list)
    (toCheck : int)
    : bool
    =
    let contains = List.contains doesContain ints
    if contains then
        incr positives
    else
        incr negatives

    create ints
    |> IntervalSet.contains toCheck
    |> (=) contains
\end{fslisting}
\end{frame}

\begin{frame}[fragile]
\frametitle{Invoke the augmented test}
\begin{fslisting}
[<Test>]
let test () =
    let pos = ref 0
    let neg = ref 0
    Check.QuickThrowOnFailure (property pos neg)

    let pos = pos.Value
    let neg = neg.Value

    pos |> shouldBeGreaterThan 0
    neg |> shouldBeGreaterThan 0

    (float pos) / (float pos + float neg)
    |> shouldBeGreaterThan 0.1
\end{fslisting}
\end{frame}

\begin{frame}
At least a tenth of the time, we want to be hitting positive cases, but \dots
\end{frame}

\begin{frame}[fragile]
\frametitle{The test is a bit unreliable!}
\begin{verbatim}
Expected: 0.1
Actual: 0.08

at FsUnitTyped.TopLevelOperators.shouldBeGreaterThan
\end{verbatim}
\end{frame}

\subsection{Manipulating the cases}
\begin{frame}
\frametitle{Manipulating the generated cases}
We want to generate cases that aren't so often ``look for something that's not in the list''.

FsCheck gives us access to its \emph{generators} for this purpose.
\end{frame}

\begin{frame}[fragile]
\frametitle{Generators}
We will have the property remain the same, but tell FsCheck to generate different cases.

FsCheck has a number of built-in generators.
It also has a computation expression to manipulate generators.
\end{frame}

\begin{frame}[fragile]
\frametitle{Generator of bounded integers}

\begin{fslisting}
let someInts : Gen<int> = Gen.choose (-100, 100)
Gen.sample 0 5 someInts
// output: [57; -24; 67; -14; 77]
\end{fslisting}
\end{frame}

\begin{frame}[fragile]
\frametitle{Generator of bounded even integers}
\begin{fslisting}
let someInts : Gen<int> = Gen.choose (-100, 100)
let evenIntegers : Gen<int> = gen {
    let! (anyInt : int) = someInt
    return 2 * anyInt
}
Gen.sample 0 5 someInts
// output: [-190; -24; -194; -108; -112]
\end{fslisting}
\end{frame}

\begin{frame}[fragile]
\frametitle{Size}
\begin{fslisting}
let integers : Gen<int list> =
    Gen.sized (fun i ->
        Gen.choose (-100, 100)
        |> Gen.listOfLength i)
\end{fslisting}
\end{frame}

\begin{frame}[fragile]
\frametitle{Generator that sometimes selects from a list}
\begin{fslisting}
let integers : Gen<int list> = ...

let listAndElt : Gen<int * int list> = gen {
    let! (list : int list) = integers
\end{fslisting}
\pause
\begin{fslisting}
    let genFromList = Gen.elements list
\end{fslisting}
\pause
\begin{fslisting}
    let genNotFromList =
        Gen.choose (-100, 100)
        |> Gen.filter (fun i -> not (List.contains i list))
\end{fslisting}
\pause
\begin{fslisting}
    let! number = Gen.oneOf [genFromList ; genNotFromList]
    return (number, list)
}
\end{fslisting}
\end{frame}

\begin{frame}[fragile]
\frametitle{Using this generator}

The generator makes pairs of an integer and a list, where the integer is 50\% likely to appear in the list.

\begin{fslisting}
[<Test>]
let test () =
    let (pos, neg) = (ref 0), (ref 0)

    (fun (list, elt) -> property pos neg elt list)
    |> Prop.forAll (Arb.fromGen listAndElt)
    |> Check.QuickThrowOnFailure

    let pos = pos.Value |> float
    let neg = neg.Value |> float

    pos / (pos + neg)
    |> shouldBeGreaterThan 0.1
\end{fslisting}
\end{frame}

\begin{frame}[fragile]
\begin{verbatim}
Ok, passed 100 tests.
\end{verbatim}

In fact we now have a positive case about 50\% of the time.
\end{frame}

\section{Stateful systems}

\begin{frame}
\tableofcontents[currentsection]
\end{frame}

\begin{frame}
\frametitle{Stateful systems}
What about state?

Key idea: \emph{describe what to do}, and then do it.
\end{frame}

\subsection{Example}
\begin{frame}[fragile]
\frametitle{Example: a stream}
Simple model: array and pointer.

\begin{table}
\captionsetup{labelformat=empty}
\caption{Starting state, pointer at index $0$}
\begin{tabular}{ |c|c|c|c|c|c|c|c }
    \hline
    \color{red}72 & 69 & 76 & 76 & 79 & 32 & 87 & \dots \\
    \hline
\end{tabular}
\end{table}

\begin{table}
\captionsetup{labelformat=empty}
\caption{Seek to index $2$}
\begin{tabular}{ |c|c|c|c|c|c|c|c }
    \hline
    72 & 69 & \color{red}76 & 76 & 79 & 32 & 87 & \dots \\
    \hline
\end{tabular}
\end{table}

\begin{table}
\captionsetup{labelformat=empty}
\caption{Write $88$ at the current index}
\begin{tabular}{ |c|c|c|c|c|c|c|c }
    \hline
    72 & 69 & \color{red}88 & 76 & 79 & 32 & 87 & \dots \\
    \hline
\end{tabular}
\end{table}

\end{frame}

\begin{frame}[fragile]
Imagine we can't access the implementation, but we want to test it anyway.

\begin{fslisting}
type Stream

[<RequireQualifiedAccess>]
module Stream =
    val uninitialised : unit -> Stream

    val read : Stream -> byte
    val write : Stream -> byte -> unit
    val seek : Stream -> int -> unit
    val currentIndex : Stream -> int
\end{fslisting}
\end{frame}

\begin{frame}
\frametitle{Some things to test}
\begin{itemize}
    \item Write then read
    \item Seek then get index
    \item Write, seek away, seek back, read
\end{itemize}

\hfill \break
FsCheck will do these, and do them well, but they are all quite specific.

Isn't the point of FsCheck to take this drudgery away from us?
\end{frame}

\begin{frame}
\frametitle{Obstacles to FsCheck}

Why is FsCheck not helping here?

\begin{itemize}
    \item Testing a mutable object
    \item No obvious immutable model to use
    \item How to generate random streams?
    \item Need shrinking not to interfere with itself
\end{itemize}
\end{frame}

\begin{frame}
Answer: \emph{describe} what to do, and \emph{then} do it!

\pause
\hfill \break

\tiny{c.f. initial algebras}

\end{frame}

\subsection{Testing with FsCheck}

\begin{frame}[fragile]
\frametitle{Domain model}
\begin{fslisting}
type StreamInteraction =
| Write of byte
| Read
| Seek of int
| CurrentIndex

type TestCase = StreamInteraction list
\end{fslisting}

\end{frame}
\begin{frame}[fragile]
\begin{fslisting}
[<RequireQualifiedAccess>]
module TestCase =
    let rec prepareStream
        (s : Stream)
        (instructions : TestCase)
        =
        for instruction in instructions do
            match instructions with
            | Write b ->
                Stream.write s b
            | Read ->
                Stream.read s |> ignore
            | Seek n ->
                Stream.seek s n
            | CurrentIndex ->
                Stream.currentIndex s |> ignore

\end{fslisting}
\end{frame}

\begin{frame}
\frametitle{Suddenly an immutable model appeared!}

\begin{itemize}
\item By constructing test cases through their \emph{descriptions}\dots
\item we made an immutable model of the world.
\item FsCheck can generate these things completely automatically!
\end{itemize}
\end{frame}

\begin{frame}[fragile]
\frametitle{Immediately useful...}
This can be used with no further modification to check a very useful property:

\begin{fslisting}
[<Test>]
let doesNotCrash () =
    let property (instructions : StreamInteraction list) =
        let s = Stream.uninitialised ()
        TestCase.prepareStream s instructions
        true

    Check.QuickThrowOnFailure (property)
\end{fslisting}
\end{frame}

\begin{frame}
\frametitle{... and more useful with generalisation}
But it really shines with just a little more work.
\end{frame}

\begin{frame}
\frametitle{The index only changes when you expect it to}
\begin{enumerate}
\item Make a generator for all interactions that shouldn't modify the index
\item Generate a list of interactions
\item Generate a list of interactions which don't modify the index
\item Concatenate
\item Execute, and assert that the index hasn't changed.
\end{enumerate}
\end{frame}

\begin{frame}
\frametitle{The gold standard: an immutable model}

\begin{enumerate}
\item Define a map of ``index'' to ``current value''
\item Interpret the interactions as acting on that map
\item Generate a list of interactions
\item Verify that the accessible output from the stream is indistinguishable from the output of the map.
\end{enumerate}
\end{frame}

\begin{frame}
\frametitle{Aside: model-based testing}

This has a name: \emph{model-based testing}.

\hfill \break

\noindent\fbox{%
Checking behaviour relative to a more easily specified model.
}

\hfill \break

MBT is a subset of property-based testing.

The property is ``the system behaves like the model does''.

(Contrast: serialiser. Correctness defined by property, not model.)

\hfill \break

FsCheck also has dedicated built-in support for MBT.
\end{frame}

\begin{frame}
\frametitle{Sketch: Testing a UI}
Test a UI by\dots
\begin{enumerate}
    \item Identifying the actions you want to test;
    \item Generating lists of actions;
    \item Applying the actions;
    \item Checking the final state is as expected given the actions.
\end{enumerate}
\end{frame}

\begin{frame}
\frametitle{Small properties, incremental addition}

The UI example is good:
\begin{itemize}
\item Enormous space of possible actions
\item Some actions very hard to test?
\item Don't care uniformly about correctness
\end{itemize}

\hfill \break

Don't need to test everything!
\end{frame}

\begin{frame}
\frametitle{A restricted search space}

If you're struggling to find properties, restrict the search space!

If you're struggling to define correctness, restrict the search space!

\hfill \break

Cut down to a small subset of user interactions, and you will find properties.
\end{frame}

\begin{frame}
\frametitle{Doge Analytics}

(Here is where I waffle, because these slides are on public Github)
\end{frame}

\section*{Summary}
\begin{frame}
\frametitle{Why you should use PBT}
Property-based testing...
\begin{itemize}
\item is easy
\item exists in many languages (Python, F\#, Haskell, \dots)
\item can be added incrementally
\item is as comprehensive as you want
\item tests anything (black-box or otherwise)
\end{itemize}
\end{frame}

\end{document}
